\section{Изложение}
\label{sec:izlozhenie}

\subsection{Конституционна и законова основа на прокурорските правомощия}
\label{sec:osnova}

Конституцията урежда прокуратурата в глава шеста „Съдебна власт“. Според чл.~127, т.~1--3
КРБ прокуратурата ръководи разследването и упражнява надзор за законосъобразното му
провеждане, може да извършва разследване, привлича към отговорност извършителите на
престъпления и поддържа обвинението по делата от общ характер. НПК конкретизира тези
функции в чл.~46: прокурорът повдига и поддържа обвинението, като ръководи разследването в
качеството на наблюдаващ прокурор, осъществява постоянен надзор за законосъобразното и
своевременното му провеждане и може сам да извършва разследване или отделни действия по
разследването.

Органи на досъдебното производство са прокурорът и разследващите органи (чл.~193 НПК),
но те не са равнопоставени. Следователите, разследващите полицаи, разследващите
митнически инспектори и разследващите инспектори от Комисията за противодействие на
корупцията действат под ръководството и надзора на прокурор (чл.~52 НПК). Самата
прокуратура е изградена йерархично. Прокурор от по-горестоящата прокуратура може да
отменя или изменя постановленията на прокурорите от по-ниската по степен прокуратура,
които не са били разгледани от съда, и да им дава задължителни писмени указания (чл.~46,
ал.~4 НПК; чл.~143 ЗСВ), а главният прокурор осъществява надзор за законност и методическо
ръководство върху дейността на всички прокурори (чл.~126, ал.~2 КРБ). Извън тази йерархия
остават европейските делегирани прокурори, когато действат по дела от компетентността на
Европейската прокуратура (чл.~46, ал.~3 и 6 НПК).\footnote{Европейската прокуратура е
създадена с Регламент (ЕС) 2017/1939 на Съвета от 12 октомври 2017~г. (ОВ L~283,
31.10.2017~г., с.~1).}

В доктрината прокурорът се разглежда като ръководител на досъдебното
производство,\footnote{Вж. Николова, Л. Прокурорът в досъдебното производство според
новия Наказателно-процесуален кодекс.~-- Адвокатски преглед, 2006, №~2--3, с.~4--22;
Чинова, М. Досъдебното производство по НПК: теория и практика. София: Сиела, 2013;
Вучков, В. Наказателен процес на Република България. София: Сиби, 2026.} а положението му
може да се опише с латинския израз \textit{dominus litis}~-- „господар на делото“.
Прокурорът решава дали производството ще започне и как ще завърши, и отговаря за
законосъобразността на действията на разследващите органи. Същевременно той е бъдещата
страна в съдебния процес, която ще поддържа обвинението въз основа на събраните под
негово ръководство доказателства. Тази двойнствена позиция обяснява защо най-тежките
намеси в правата на гражданите са поставени под предварителен съдебен контрол
(вж. част~\ref{sec:prinuda}).

\subsection{Предварителна проверка и образуване на досъдебното производство}
\label{sec:obrazuvane}

Досъдебно производство се образува, когато са налице законен повод и достатъчно данни за
извършено престъпление (чл.~207, ал.~1 НПК). Законните поводи са изчерпателно изброени в
чл.~208 НПК, а анонимното съобщение не е законен повод (чл.~209, ал.~1 НПК). Достатъчни
данни са налице, когато може да се направи основателно предположение, че е извършено
престъпление, без да е нужно извършителят да е известен (чл.~211 НПК). Това е съществено
при киберпрестъпленията, при които пострадалият обикновено знае какво е загубил, но не и
кой стои зад атаката. Производството може да започне срещу неизвестен извършител, а
самоличността му да бъде установена по-късно, например чрез данните за IP адресите.

Когато сведенията не са достатъчни, прокурорът може да извърши предварителна проверка:
да изисква документи и сведения, лично да извършва проверки или да ги възлага на други
органи. Срокът на проверката е два месеца и може да бъде удължен еднократно с един месец
(чл.~145, ал.~1 и 2 ЗСВ).\footnote{Вж. Вучков, В. Образуване на досъдебно производство: актуални
проблеми. София: Сиби, 2015; за критичен анализ~-- Иванова, И. Конституционни проблеми във
връзка с правомощията на прокурора по чл.~145 от Закона за съдебната власт.~-- Юридически
свят, 2023, №~2, с.~32--59.} В хода на проверката обаче не могат да се извършват действия по
разследването. Данните за трафика могат да бъдат поискани по чл.~159а НПК само в
образувано досъдебно производство, а прокуратурата не е сред органите с право на справка
по чл.~251в, ал.~1 ЗЕС. Тъй като доставчиците съхраняват тези данни шест месеца
(чл.~251б, ал.~1 ЗЕС), продължителната проверка може да доведе до загубата им. Затова при
данни за киберпрестъпление е за предпочитане производството да бъде образувано
незабавно: с постановление на прокурора или със съставянето на протокола за първото
неотложно действие по разследването, за което разследващият орган уведомява прокурора в
срок до 24~часа (чл.~212 НПК).

Ако няма основание за образуване на производство, прокурорът отказва да го образува
(чл.~213, ал.~1 НПК). Дълго време отказът подлежеше само на контрол от по-горестоящата
прокуратура. От 2023~г. потвърденият отказ за тежко престъпление и за изрично изброени
други престъпления подлежи на обжалване пред първоинстанционния съд, който може да го
отмени и да върне преписката със задължителни указания (чл.~213, ал.~4--6
НПК).\footnote{Разпоредбите са създадени с изменението, обн. ДВ, бр.~48 от 2023~г., с което
беше уреден и особен ред за разследване срещу главния прокурор (чл.~411а--411и НПК). Още в
решението по делото „Колеви срещу България“ от 2009~г. ЕСПЧ прие, че при тогавашната
уредба независимото разследване срещу главния прокурор е практически невъзможно. С решение
№~14 от 26.07.2024~г. по к.д. №~10/2023~г. Конституционният съд отхвърли искането
чл.~213, ал.~4--6 и чл.~411а--411и НПК да бъдат обявени за противоконституционни.} Така се
стесни кръгът на прокурорските решения, които остават извън съдебен контрол. Промяната
засяга пряко и киберпрестъпленията. След изменението на НК, обн. ДВ, бр.~53 от 2022~г.,
основните състави на компютърните престъпления по чл.~319а--319е НК се наказват с
лишаване от свобода с горна граница поне шест години, т.е. те са тежки престъпления по
смисъла на чл.~93, т.~7 НК. Тежко престъпление е и компютърната измама по чл.~212а НК.

\subsection{Ръководство и надзор върху разследването}
\label{sec:rakovodstvo}

Когато не разследва сам, прокурорът упражнява ръководство и надзор. Той може да
контролира непрекъснато хода на разследването, да дава указания, да участва в действията
по разследването или да ги извършва сам, да отстранява разследващия орган, да изземва
делото и да го предава на друг орган, да възлага отделни действия на органите на МВР,
ДАНС, Агенция „Митници“ и Комисията за противодействие на корупцията и да отменя
постановленията на разследващия орган (чл.~196, ал.~1 НПК). Писмените му указания са
задължителни и не подлежат на възражение (чл.~197 НПК). Когато се съберат достатъчно
доказателства за виновността на определено лице, разследващият орган докладва на
прокурора и привлича лицето като обвиняем (чл.~219, ал.~1 НПК). Така прокурорът
контролира обвинението още от формулирането му.\footnote{Вж. Манев, А. Ролята на
прокурора при привличането на обвиняемия и повдигането на предварителното обвинение.~--
Общество и право, 2023, №~8, с.~44--60.}

Наблюдаващият прокурор следи и за срока на разследването (чл.~196, ал.~2 НПК). След
изменението, обн. ДВ, бр.~65 от 2025~г., разследването трябва да приключи в три месеца, а
при фактическа и правна сложност~-- в шест месеца от образуването на производството. В
такива случаи прокурорът може да удължи срока до дванадесет месеца, а следващите
удължавания се разрешават от административния ръководител на прокуратурата (чл.~234,
ал.~1--3 НПК).

При киберпрестъпленията ръководството на разследването има и организационна страна.
Делата с фактическа и правна сложност могат да бъдат възложени на следовател (чл.~194,
ал.~1, т.~4 НПК), а отделни действия~-- на специализираните звена на МВР. Обикновено е
необходима и компютърно-техническа експертиза (чл.~144, ал.~1 НПК). От точността на
въпросите към експерта зависи дали ще бъдат установени релевантните факти: кой е имал
достъп до системата, кога и от кой адрес, кои данни са били променени и как са били
прехвърлени средствата.

\subsection{Прокурорът и съдебният контрол върху процесуалната принуда}
\label{sec:prinuda}

Законът разграничава правомощията, които прокурорът упражнява сам, от тези, при които той
само прави искане, а решение взема съдът \tabref{tab:kontrol}. Колкото по-тежко е
засягането на основните права (личната свобода, неприкосновеността на жилището, тайната
на кореспонденцията), толкова по-необходим е предварителният съдебен контрол.

\begin{table}[t]
\centering
\caption{Правомощия на прокурора и съдебен контрол в досъдебното производство}
\label{tab:kontrol}
\small
\begin{tabularx}{\textwidth}{@{}L{3.9cm}YL{3.4cm}L{2.5cm}@{}}
\toprule
\textbf{Мярка или действие} & \textbf{Роля на прокурора} & \textbf{Кой решава} & \textbf{Основание}\\
\midrule
Задържане под стража; домашен арест & Прави искане; може да задържи обвиняемия до 72~часа за довеждането му пред съда & Първоинстанционният съд & чл.~62, ал.~2; чл.~64 НПК\\
Подписка; гаранция & Взема мярката, както и разследващият орган & Органът на досъдебното производство; гаранцията се обжалва пред съда & чл.~58; чл.~61, ал.~3 НПК\\
Забрана за напускане на страната & Налага забраната и разрешава пътуванията & Прокурорът; отказът се обжалва пред съда & чл.~68 НПК\\
Претърсване и изземване, вкл. на компютърни данни & Прави искане; при неотложност представя протокола за одобряване до 24~часа & Съдия от първоинстанционния съд & чл.~160, 161 НПК\\
Задържане и изземване на кореспонденция и електронна поща & Прави искане (само при тежки престъпления) & Съдия от първоинстанционния съд & чл.~165 НПК\\
Данни за трафика от доставчици на електронни съобщения & Прави мотивирано искане за период до 6~месеца & Съдия от първоинстанционния съд & чл.~159а НПК\\
Специални разузнавателни средства & Прави писмено мотивирано искане & Председателят на окръжния съд & чл.~173, 174 НПК\\
Запазване и предаване на компютърни данни & Изисква ги пряко от лицата, у които се намират & Прокурорът & чл.~159 НПК\\
Отваряне на носител с иззети компютърни данни & Дава разрешение & Прокурорът & чл.~163, ал.~16 НПК\\
Криптоактиви & Иззетите активи се прехвърлят в портфейл, контролиран от прокуратурата; прокурорът може да разреши продажбата им & Прокурорът & чл.~163, ал.~15 и 17 НПК\\
\bottomrule
\end{tabularx}
\end{table}

Мярката задържане под стража се взема от съда по искане на прокурора, който може само
да задържи обвиняемия до 72~часа за довеждането му пред съда (чл.~64, ал.~1 и 2 НПК). В
неотложни случаи претърсване и изземване може да се извърши и без разрешение, но
наблюдаващият прокурор представя протокола за одобряване от съдията в срок до 24~часа
(чл.~161, ал.~2 НПК). Във всички случаи, в които решава съдът, прокурорът е филтър:
преценява дали мярката е необходима и съразмерна, мотивира искането и отговаря
за законосъобразното ѝ изпълнение, а съдът проверява дали са налице законовите предпоставки.

\subsection{Приключване на досъдебното производство}
\label{sec:priklyuchvane}

След приключването на разследването разследващият орган изпраща делото на прокурора с
писмено мнение (чл.~235, ал.~1 НПК). В срок до един месец, а при фактическа и правна
сложност~-- до два месеца, прокурорът прекратява или спира производството, внася в съда
споразумение или предложение за освобождаване от наказателна отговорност с налагане на
административно наказание, или повдига обвинение с обвинителен акт (чл.~242, ал.~1, 4 и
5 НПК). Ако е нужно, той отстранява допуснатите процесуални нарушения или извършва
допълнителни действия по разследването (чл.~242, ал.~2 и 3 НПК).\footnote{Вж. Вучков, В.
Наказателен процес на Република България. Цит. съч.; Чинова, М., Г.~Митов. Кратък
лекционен курс по наказателнопроцесуално право. София: Сиела, 2021; Манев, Н.
Наказателнопроцесуално право: ръководство за студенти. София: Сиела, 2016.} Изборът между
тези възможности е най-яркият израз на ролята на прокурора като \textit{dominus litis}: от
неговото решение зависи дали делото ще стигне до съда.

Прокурорът съставя обвинителен акт, когато е убеден, че са събрани необходимите
доказателства (чл.~246, ал.~1 НПК). В акта се описват престъплението, времето, мястото и
начинът на извършването му (чл.~246, ал.~2 НПК). При киберпрестъпленията начинът на
извършване включва използвания софтуер, мрежовите адреси, времето на достъп и движението
на средствата, а неясното им описание затруднява защитата и излага обвинението на
оспорване.

Прекратяването и спирането подлежат на съдебен контрол по жалба на обвиняемия,
пострадалия или ощетеното юридическо лице (чл.~243, ал.~4--6 и чл.~244, ал.~3 и 5 НПК).
Необжалваното постановление за прекратяване може да бъде отменено от по-горестоящата
прокуратура в срок до две години при тежко престъпление и до една година в останалите
случаи. След изменението, обн. ДВ, бр.~97 от 2025~г., главният прокурор може да го отмени
и след този срок на изчерпателно изброени основания (чл.~243, ал.~10 и 11 НПК).
Производството за киберпрестъпление може да бъде спряно, когато извършителят не е разкрит
(чл.~244, ал.~1, т.~2 НПК). Тогава прокурорът изпраща делото на съответните органи за
издирване на извършителя и възобновява производството, когато основанието за спиране
отпадне (чл.~245 НПК). Споразумението може да служи и за обезщетяване на пострадалите,
защото при причинени имуществени вреди то се допуска само след възстановяването или
обезпечаването им (чл.~381, ал.~3 НПК).

\subsection{Особености на ролята на прокурора при киберпрестъпленията}
\label{sec:kiber}

Киберпрестъпленията в тесен смисъл са уредени в глава девета „а“ на Особената част на
НК~-- „Компютърни престъпления“ (чл.~319а--319е), а в по-широк смисъл към тях се отнасят и
престъпленията, извършени чрез информационна система, например компютърната измама по
чл.~212а НК.\footnote{За разследването на компютърни престъпления вж. Владова-Недкова, И.
Разследване на компютърни престъпления. София: Софи-Р, 2023.} Доказателствата за тях са
електронни и лесно се променят, съхраняват се от частни доставчици за ограничено време и
често се намират в друга държава, а самоличността на дееца обикновено е скрита. Затова
бързината и законосъобразността на действията на прокурора решават дали доказателствата
ще бъдат запазени и използвани пред съда.

\subsubsection{Запазване и събиране на електронни доказателства}

Първата задача на прокурора е да осигури запазването на данните. Чл.~159 НПК задължава
учрежденията, юридическите лица, длъжностните лица и гражданите да запазят и да предадат
по искане на органите на досъдебното производство намиращите се у тях компютърни
информационни данни, които могат да имат значение за делото.\footnote{За електронните
доказателства в българския наказателен процес вж. Вучков, В. Електронните доказателства
по наказателни дела и техните източници~-- необходими разсъждения.~-- Юридическо
списание на НБУ, 2023, т.~19, №~2, с.~12--19; Младенов, Д. Електронни доказателства.~--
De Jure, 2022, т.~13, №~1, с.~63--74.} Разпоредбата отговаря на изискването на
Конвенцията за престъпления в кибернетичното пространство за бързо запазване на
съхранени компютърни данни (чл.~16).\footnote{Конвенция за престъпления в кибернетичното
пространство, ратифицирана със закон, обн. ДВ, бр.~29 от 2005~г.; в сила за Република
България от 01.08.2005~г. Вж. и Explanatory Report to the Convention on Cybercrime.
European Treaty Series No.~185. Budapest, 2001.}

Достъпът до данните за трафика е уреден по-строго. Доставчиците на обществени електронни
съобщителни мрежи и услуги ги предоставят в досъдебното производство само въз основа на
мотивирано разпореждане на съдия, издадено по искане на наблюдаващия прокурор, и само при
разследването на тежки умишлени престъпления, а периодът на справката не може да
надхвърля шест месеца (чл.~159а, ал.~1, 2 и 5 НПК). Разпоредбата е създадена през 2015~г.,
след като Конституционният съд обяви за противоконституционна предишната уредба на
съхраняването на данни и на достъпа до тях.\footnote{РКС №~2 от 12.03.2015~г. по к.д.
№~8/2014~г., обн. ДВ, бр.~23 от 2015~г. Чл.~159а НПК е създаден с изменението, обн. ДВ,
бр.~24 от 2015~г.} Съдът на ЕС прие, че прокуратура, която ръководи досъдебното
производство и впоследствие поддържа обвинението, не може да разрешава достъп до данни за
трафика и за местоположението.\footnote{Решение на СЕС от 2 март 2021~г. по дело C-746/18,
\textit{Prokuratuur}, ECLI:EU:C:2021:152.} По запитване на български съд той прие и че
правото на ЕС не допуска общо и неизбирателно превантивно съхраняване на такива данни,
дори когато то е ограничено до шест месеца.\footnote{Решение на СЕС от 17 ноември 2022~г.
по дело C-350/21, \textit{Spetsializirana prokuratura}, ECLI:EU:C:2022:896.} Това поставя
под въпрос шестмесечното съхраняване по чл.~251б ЗЕС. ЕСПЧ на свой ред установи, че
българската уредба на съхраняването на комуникационни данни и на достъпа до тях не
съдържа достатъчни гаранции срещу злоупотреби.\footnote{Решение на ЕСПЧ от 11 януари
2022~г. по делото „Екимджиев и други срещу България“, жалба №~70078/12.} Затова моделът,
при който прокурорът иска, а съдът разрешава, е минималният стандарт, а качеството на
мотивите в искането на прокурора е гаранция за законосъобразността на доказателствата.

При претърсване прокурорът иска разрешение, което обхваща и компютърните информационни
системи (чл.~160, ал.~1 НПК). Иззетите компютърни данни се записват върху носител, който
се запечатва и се пази от физическо и електромагнитно въздействие. Отварянето му за
нуждите на разследването става само с разрешение на прокурора, в присъствието на поемни
лица и на специалист~-- технически помощник (чл.~163, ал.~14 и 16
НПК).\footnote{Вж. Владова-Недкова, И. Цит. съч.} Така прокурорът става гарант за
целостта на доказателствата, която на практика се удостоверява с контролни суми (хеш
стойности), изчислени при изземването и проверявани при всяко отваряне на
носителя.\footnote{Вж. Council of Europe. Electronic Evidence Guide. Version~2.1.
Strasbourg: Council of Europe, 2020.} Аналогичен стандарт Съдът на ЕС прилага към
мобилните телефони: достъпът на полицията до данните в тях по правило изисква
предварителен контрол от съд или независим орган, освен в надлежно обосновани спешни
случаи.\footnote{Решение на СЕС от 4 октомври 2024~г. по дело C-548/21,
\textit{Bezirkshauptmannschaft Landeck}, ECLI:EU:C:2024:830.} Затова е препоръчително
искането на прокурора за претърсване и изземване да обхваща изрично и прегледа на данните
в иззетите устройства.

С изменението на НПК, обн. ДВ, бр.~65 от 2025~г., беше уредено и изземването на
криптоактиви. Те се прехвърлят на адрес на портфейл или на друга платформа, контролирани
от Прокуратурата на Република България, а по искане на лицето, от което са иззети,
прокурорът може да разреши продажбата им (чл.~163, ал.~15 и 17 НПК). Правилото отчита,
че криптоактивите не могат да бъдат запечатани като вещ, защото всеки, който знае
частния ключ, може да ги прехвърли.

\subsubsection{Международно сътрудничество}

Основният международен договор в тази област е Конвенцията за престъпления в
кибернетичното пространство. Тя урежда бързото запазване на данни по молба на друга
държава (чл.~29) и мрежа от звена за контакт, достъпни 24 часа в денонощието (чл.~35).
Централен орган за молбите за правна помощ в досъдебното производство е Върховната
касационна прокуратура, а звеното за контакт е в Главна дирекция „Борба с организираната
престъпност“ на МВР.\footnote{Декларации на Република България по чл.~27 и 35 от
Конвенцията, публикувани в базата данни на договорите на Съвета на Европа. Вторият
допълнителен протокол към Конвенцията (CETS No.~224) е подписан от България на
12.05.2022~г., но към ноември 2026~г. не е ратифициран.} Върховната касационна
прокуратура създава и съвместни екипи за разследване с други държави (чл.~476, ал.~3 НПК).

В рамките на ЕС прокурорът издава европейска заповед за разследване. Когато обаче
българският закон изисква за съответното действие разрешение на съдия, заповедта се
издава след такова разрешение или от съда по искане на прокурора (чл.~5, ал.~1, т.~1
ЗЕЗР). Уредбата съответства на решението на СЕС по дело C-724/19 (отново по
запитване на български съд), според което прокурорът не може да издаде такава заповед за данни за
трафика и за местоположението, когато в сходен национален случай достъпът до тях се
разпорежда само от съдия.\footnote{Решение на СЕС от 16 декември 2021~г. по дело C-724/19,
ECLI:EU:C:2021:1020. Чл.~5, ал.~1 ЗЕЗР е изменен с ДВ, бр.~36 от 2024~г.}

От 18~август 2026~г. се прилага Регламент (ЕС) 2023/1543 (чл.~34, ал.~2). Той позволява
на органа от една държава членка да се обърне директно към доставчик на услуги в друга
държава членка с европейска заповед за предоставяне или за запазване на електронни
доказателства. Заповед за данни за абонат и за данни, поискани единствено за
идентифициране на ползвателя, както и заповед за запазване на данни, може да издаде и
прокурор. Заповедта за предоставяне на данни за трафика и на данни за съдържанието обаче
се издава или потвърждава само от съдия или съд (чл.~4). Освен това заповедта е
допустима само ако подобна заповед би могла да бъде издадена в сходен национален случай
(чл.~5, ал.~2), поради което българският прокурор трябва да се съобразява и с изискването
на чл.~159а НПК за съдебно разрешение. Доставчикът предава данните в срок до десет дни, а в
спешни случаи~-- до осем часа (чл.~10). Към
ноември 2026~г. обаче в България не са приети национални правила за прилагането на
регламента, а срокът за транспониране на придружаващата го Директива (ЕС) 2023/1544
изтече на 18~февруари 2026~г. Проект за изменение на НПК, който урежда тези въпроси, беше публикуван за обществено обсъждане през септември
2026~г.\footnote{Обществена консултация №~12669-K на Министерството на правосъдието на
Портала за обществени консултации (strategy.bg), 03.09.--05.10.2026~г.}
